%%%%%%%%%%  scan idx 509  |  folio 496  %%%%%%%%%%
\SGAPrintPage{496}%
Let us now calculate the second side of the Riemann--Roch formula. For every \(k\), we have:
\[
 \widetilde c(\xi^k)=(\widetilde c(\xi))^k=(1,1+x)^k=(1,1+kx)
\]
by (V (6.2.1)). Consequently:
\sgaStableTarget{sga6.e08.eq.5-4-2}\begin{equation}\tag{5.4.2}
 \ch(\xi^k)=\exp(kx).
\end{equation}
The second side of the Riemann--Roch formula is therefore:
\sgaStableTarget{sga6.e08.eq.5-4-3}\begin{equation}\tag{5.4.3}
 f_{\mathrm{gr}}\left(\exp(kx)\cdot
 \frac{x^{r+1}}{(1-\exp(-x))^{r+1}}\right)
\end{equation}
Since \(x\) is the class in \(\Gr^1K^{\bullet}(P)\) of \(1-\xi^{-1}\), \(f_{\mathrm{gr}}(x^n)\) is, for \(0\leq n\leq r\), the class in
\(\Gr^{n-r}K^{\bullet}(Y)\) of \(f_*((1-\xi^{-1})^n)=1\). We thus have:
\sgaStableTarget{sga6.e08.eq.5-4-4}\begin{equation}\tag{5.4.4}
 f_{\mathrm{gr}}(x^n)=0\quad\text{if }0\leq n<r~;\qquad
 f_{\mathrm{gr}}(x^r)=1.
\end{equation}
Finally recall that the elements \(1,x,\ldots,x^{r+1}\) are linearly dependent over \(\Gr K^{\bullet}(Y)\) (VI 5.6), and satisfy the equation:
\sgaStableTarget{sga6.e08.eq.5-4-5}\begin{equation}\tag{5.4.5}
 x^{r+1}=0,
\end{equation}
and that \(1,\ldots,x^r\) are linearly independent over \(\Gr K^{\bullet}(Y)\).

One can then write the series \(\varphi_k=\exp(kx)\cdot\frac{x^{r+1}}{(1-\exp(-x))^{r+1}}\) uniquely as a linear combination, with coefficients in
\(\widehat{\Gr K^{\bullet}(Y)_{\mathbb Q}}\), of the elements \(1,\ldots,x^r\). The image of \(\varphi_k\) under \(f_{\mathrm{gr}}\) is then, by (5.4.4) and 3.7, its coefficient of \(x^r\). Taking (5.4.1) into account, the Riemann--Roch theorem for \(f\) is therefore equivalent to:
\sgaStableTarget{sga6.e08.eq.rr}\[
\left\{
\begin{array}{ll}
\text{for }-r\leq k<0,\text{ the coefficient of }x^r\text{ in }\varphi_k\text{ is }0~;\\
\text{for }k=0,\text{ the coefficient of }x^r\text{ in }\varphi_k\text{ is }1.
\end{array}
\right.
\tag{RR}
\]
This coefficient is the residue of the differential form
\[
 \frac{\exp(kx)\cdot dx}{(1-\exp(-x))^{r+1}}.
\]

%%%%%%%%%%  scan idx 510  |  folio 497  %%%%%%%%%%
\SGAPrintPage{497}%
Putting
\(y=1-\exp(-x)\), we are reduced to calculating the residue of:
\[
 \frac{(1-y)^{-k-1}\cdot dy}{y^{r+1}},
\]
that is, the coefficient of \(y^r\) in \((1-y)^{-k-1}\). For \(-r\leq k<0\), it is zero, and for \(k=0\), it is \(1\).

This completes the proof of the Riemann--Roch theorem.

\sgaStableTarget{sga6.e08.bibliography}\subsection*{BIBLIOGRAPHY}
\begin{enumerate}[label={[\arabic*]}]
\item A. Borel and J.P. Serre, Le théorème de Riemann--Roch, Bull. Soc. Math. de France, t. 86, 1958.
\item A. Grothendieck, Techniques de construction et théorèmes d'existence en Géométrie Algébrique, Sém. Bourbaki, 1960-61.
\end{enumerate}
